\chapter{The envelope theorem: statement and sharpness}
\index{envelope theorem!unibranch projective variety}\index{envelope property!sharpness}
\label{chap:bj-envelope-statement}

This chapter states the envelope theorem with every hypothesis visible and
gives its proof immediately after the statement.
\Crefrange{chap:bj-cartier-realization}{chap:bj-subintegral-normalization}
establishes, in the
order in which they are needed, Cartier realization, two reflection
theorems, smooth regularization, finite descent, alterations, and
normalization reflection.  The proof below invokes those results in that
order; \Cref{chap:bj-envelope-consequences} records a second assembly and
its consequences.

Throughout, $k$ is algebraically closed and carries the trivial absolute
value.  The scheme $X$ is an integral projective $k$-variety, and
$\theta\in N^1(X)_{\mathbf R}$.  We use the Boucksom--Jonsson class
$\PSH_X(\theta)$ fixed in
\Crefrange{chap:bj-fs-pl-psh}{chap:bj-envelope-history}.  Its members take values in
$\mathbf R\cup\{-\infty\}$: they are finite at divisorial valuations but
may equal $-\infty$ at semivaluations with proper support.


\section{Orientation}
\label{sec:ch37-orientation}

\subsection{The theorem and its place in the book}

This chapter states the principal theorem of Part~IV and gives its proof.
The final assembly is conceptually short: it reduces the general case to the
normal case by normalization, while every input it uses is established in
\Crefrange{chap:bj-cartier-realization}{chap:bj-subintegral-normalization}.
The reader who wants the mathematics rather than the bookkeeping should read
the statement, then \Cref{sec:ch37-proof-map}, and then the chapters in
order.

What makes the result worth nine chapters is not the final assembly but the
descent.  Two independent difficulties have to be overcome, and they are
solved by unrelated techniques:

\begin{itemize}
\item \emph{the smooth case in arbitrary characteristic}, where multiplier
      ideals are unavailable in characteristic $p$ and are replaced by
      asymptotic test ideals, with the subadditivity and uniform
      $O(m^{-1})$ control that this substitution demands
      (\Cref{chap:bj-smooth-characteristic});
\item \emph{descent along the normalization}, which is finite and birational
      but not flat, not separable in general, and not an isomorphism.  In
      characteristic $p$ it is handled by Frobenius and universal
      homeomorphisms (\Cref{chap:bj-universal-homeomorphisms}); in
      characteristic zero by the Roberts--Singh additive-power sandwich for
      subintegral extensions (\Cref{chap:bj-subintegral-normalization}).
\end{itemize}

\subsection{Why the hypotheses are what they are}

The statement mixes three different kinds of hypothesis, and keeping their
roles separate prevents false sharpness claims:

\begin{itemize}
\item \emph{Unibranchness} is the geometric condition.  It is necessary
      when the class is ample, by \Cref{thm:ep-implies-unibranch}; the
      node--cusp comparison of \Cref{sec:ch37-node-cusp} diagnoses the
      normalization-reflection obstruction but is not itself the converse
      theorem.
\item \emph{Nonemptiness and a common upper bound} are admissibility
      conditions on the family.  Nonemptiness excludes the formal supremum
      $-\infty$, which is not psh under \Cref{def:bj-psh}\,(2), while
      \Cref{ex:ch37-pointwise-not-uniform} shows that the common upper bound
      cannot be dropped.  For a family of psh functions, being bounded at
      every analytic point is already equivalent to having one common
      constant bound; see \Cref{rem:ch37-pointwise-equals-uniform}.
\item \emph{Integrality, projectivity, algebraic closedness, and the trivial
      absolute value} delimit the theory proved in Part~IV.  They are not
      presented here as a sharp list for every possible extension of the
      theorem.  In particular, projectivity supplies compactness of
      $X^{\An}$ and the numerical and section-theoretic apparatus used by
      the proof.
\end{itemize}

Nothing is assumed about the characteristic, the dimension, the rationality
of $\theta$, or the behaviour of the $\varphi_i$ at proper-support
semivaluations.  Each of these was a genuine restriction in some earlier
partial result, and \Cref{sec:ch37-hypothesis-ledger} records which.

\section{The exact statement}

\begin{definition}[Bounded-above family]
\label{def:bj-bounded-family}
A family
\[
              (\varphi_i)_{i\in I}\subset\PSH_X(\theta)
\]
is \emph{bounded above} if there exists $C\in\mathbf R$ such that
$\varphi_i\le C$ on $X^{\An}$ for every $i\in I$.
\end{definition}

\begin{remark}[Pointwise and uniform bounds agree here]
\label{rem:ch37-pointwise-equals-uniform}
For families inside $\PSH_X(\theta)$ the three conditions
\[
 \sup_i\sup_{X^{\An}}\varphi_i<\infty,
 \qquad
 \sup_i\varphi_i(v)<\infty\ \text{ for every }v,
 \qquad
 \sup_i\varphi_i(v_{\triv})<\infty
\]
are equivalent, so \Cref{def:bj-bounded-family} is not strengthening a
pointwise hypothesis into a uniform one.  The reason is the maximum
principle at the trivial valuation: every Fubini--Study function satisfies
$\sup_{X^{\An}}\phi=\phi(v_{\triv})$ by
\Cref{prop:ch34-generic-finiteness}\,\cref{itm:ch34-gf:decreasing}, and if
$\phi_j\searrow\varphi$ is an approximating net as in
\Cref{def:bj-psh}, then
$\sup_{X^{\An}}\varphi\leq\inf_j\sup_{X^{\An}}\phi_j
 =\inf_j\phi_j(v_{\triv})=\varphi(v_{\triv})\leq\sup_{X^{\An}}\varphi$,
so $\sup_{X^{\An}}\varphi=\varphi(v_{\triv})$ for every
$\varphi\in\PSH_X(\theta)$.  Hence a bound at the single point
$v_{\triv}$ is already a bound on $X^{\An}$.  What \emph{does} fail is
any attempt to drop the bound: see
\Cref{ex:ch37-pointwise-not-uniform}.
\end{remark}

\begin{definition}[Envelope property]
\label{def:bj-envelope-property-final}
For a nonempty bounded-above family define
\[
 u_0:=\sup_{i\in I}\varphi_i,
 \qquad
 P(\varphi_i):=u_0^*,
\tag{37.1}\label{eq:bj-envelope-definition}
\]
where $u_0^*$ is the least upper semicontinuous majorant of $u_0$.
We write $\operatorname{EP}_X(\theta)$ if
\[
                    P(\varphi_i)\in\PSH_X(\theta)
\]
for every such family.
\end{definition}

Nonemptiness excludes the formal supremum $-\infty$, which is not a psh
function.  The star cannot be suppressed: an arbitrary supremum of upper
semicontinuous functions need not be upper semicontinuous.

\begin{theorem}\label{thm:bj-final-envelope}
Let $X$ be an integral projective unibranch variety over an algebraically
closed trivially valued field $k$, and let
$\theta\in N^1(X)_{\mathbf R}$.  If
\[
               (\varphi_i)_{i\in I}\subset\PSH_X(\theta)
\]
is nonempty and bounded above, then
\[
             \left(\sup_{i\in I}\varphi_i\right)^*
                    \in\PSH_X(\theta).
\tag{37.2}\label{eq:bj-final-envelope}
\]
The assertion holds in arbitrary dimension and characteristic, for every
real numerical class, and without a finiteness hypothesis at
proper-support semivaluations.
\end{theorem}

\begin{proof}
Let $\nu:X^\nu\to X$ be the normalization.  Since $X$ is
unibranch, \Cref{lem:bj-unibranch-normalization-universal} shows that
$\nu$ is a finite birational universal homeomorphism.  In particular,
$\nu^{\An}$ is a homeomorphism.  The normal projective variety $X^\nu$
has the envelope property in every real numerical class by
\Cref{cor:bj-normalization-source-ep}.

Pullback therefore gives a bounded-above nonempty family
$\nu^*\varphi_i\in\PSH_{X^\nu}(\nu^*\theta)$, and hence
\[
 v:=\left(\sup_i\nu^*\varphi_i\right)^*
       \in\PSH_{X^\nu}(\nu^*\theta).
\]
Because $\nu^{\An}$ is a homeomorphism, upper semicontinuous
regularization commutes with pullback by
\Cref{cor:bj-universal-usc}.  Thus, as an identity of extended-real
functions on the whole analytification,
\[
 v
 =\nu^*\left(\sup_i\varphi_i\right)^*.
\]
This remains valid at proper-support points where both sides may be
$-\infty$: it follows by transporting neighborhood filters through the
homeomorphism, without subtracting extended-real values.  Finally,
\Cref{cor:bj-universal-reflection-all-characteristics} reflects
$\nu^*$-psh functions through a finite birational universal
homeomorphism.  It follows that
$\left(\sup_i\varphi_i\right)^*\in\PSH_X(\theta)$, as required.
\end{proof}

\begin{remark}[Logical scope]
\label{rem:ch37-logical-scope}
Normality is not assumed.  In particular, the theorem includes unibranch
nonnormal varieties such as a cuspidal curve.  Conversely, this chapter
does not assert that every non-unibranch variety fails every envelope
statement.  The sharpness discussion below concerns the normalization
reflection required by this proof.
\end{remark}

\begin{remark}[Empty psh classes]
\label{rem:ch37-empty-psh-classes}
If $\PSH_X(\theta)=\varnothing$, then
$\operatorname{EP}_X(\theta)$ is vacuous.  The family in
\Cref{thm:bj-final-envelope} makes the theorem's instance nonvacuous.
\end{remark}

\section{Equivalent formulations}

Write $X^{\divis}$ for the trivial valuation together with the positive
rational multiples of divisorial valuations.  All these points have full
support.

We use exactly the comparison interface recorded in the supplied envelope
manuscript from \cite[Theorem~4.22]{BJ22}: if
$u\in\PSH_X(\theta)$ and
\[
              h:X^{\An}\longrightarrow\mathbf R\cup\{-\infty\}
\]
is upper semicontinuous, then $u\le h$ on $X^{\An}$ if and only if
the inequality holds on $X^{\divis}$.  In particular, two psh functions
that agree on $X^{\divis}$ agree everywhere.

\begin{proposition}\label{prop:bj-increasing-net-criterion}
\label{prop:bj-ep-equivalences}
For fixed $X$ and $\theta$, the following assertions are equivalent.
\begin{enumerate}
\item $\operatorname{EP}_X(\theta)$ holds.
\item Every bounded-above upward-directed net
\[
                  (u_j)_{j\in J}\subset\PSH_X(\theta)
\]
has a pointwise supremum on $X^{\divis}$ that is the divisorial trace of
a unique $u\in\PSH_X(\theta)$; moreover,
\[
                         u=\left(\sup_j u_j\right)^*.
\]
\item For every nonempty bounded-above family
$(\varphi_i)_{i\in I}$, the net
\[
                 u_F:=\max_{i\in F}\varphi_i
\tag{37.3}\label{eq:bj-finite-max-net}
\]
indexed by the finite nonempty subsets $F\subset I$ has the limit
described in the preceding item.
\end{enumerate}
\par
\end{proposition}

\begin{proof}
The equivalence of the first two assertions is
\cite[Theorem~5.11]{BJ22}.  Finite maxima in one psh class are psh.
The finite nonempty subsets of $I$, ordered by inclusion, form a directed
set, and
\[
                        \sup_F u_F=\sup_i\varphi_i.
\]
Thus the second assertion implies the third, and the third implies the
first.

Conversely, if $(u_j)$ is upward directed, every finite maximum of its
members is dominated by one member of the net.  Its finite-maximum net
therefore has the same pointwise supremum.  The third assertion then yields
the second.  Uniqueness follows from divisorial determination.
\end{proof}

\begin{corollary}\label{cor:bj-historical-conjecture}
\Cref{thm:bj-final-envelope}, applied separately to every
$\theta\in N^1(X)_{\mathbf R}$, is precisely Boucksom--Jonsson
Conjecture~5.14 for projective unibranch varieties.
\end{corollary}

\begin{proof}
Conjecture~5.14 is the assertion of the envelope property in every numerical
class.  The equivalence between its family and increasing-net forms is
\Cref{prop:bj-increasing-net-criterion}.
\end{proof}

\section{The extended-real boundary}

The theorem is an assertion on all of $X^{\An}$, not only on valuations
of $k(X)$.  Let $x\in X^{\An}$ have proper support $Z\subsetneq X$.
If a coherent ideal $\mathfrak a$ vanishes identically along $Z$, then
\[
                  v_x(\mathfrak a)=+\infty,
 \qquad
                  \log|\mathfrak a|(x)=-\infty.
\tag{37.4}\label{eq:bj-supported-ideal-value}
\]
Hence a generically finite Fubini--Study function can have value
$-\infty$ at $x$.

\begin{example}[A common base locus]
\label{ex:bj-common-base-locus}
Let
\[
 \phi=\frac1m\max_j\{\log|s_j|+\lambda_j\}
\]
be a rational GF function.  If all $s_j$ vanish along a proper closed
set $Z$, then $\phi=-\infty$ at the trivial semivaluation supported on
$Z$.  This is not a value to be discarded: determinant norms and
normalization descent must reproduce it exactly.
\end{example}

\par\medskip

\begin{proposition}\label{prop:bj-safe-interfaces}
The proof of \Cref{thm:bj-final-envelope} uses the following interfaces.
\begin{enumerate}
\item Equalities of ideal values are read before changing signs, so the
value $+\infty$ is allowed.
\item Finite maxima and finite positively weighted sums are read in
$\mathbf R\cup\{-\infty\}$.
\item Divisor identities at proper-support points are global psh identities,
not rearranged differences of two functions equal to $-\infty$.
\item A decreasing psh limit is used only after a finite value at the
trivial valuation, or an explicit GF lower bound, excludes the identically
$-\infty$ alternative.
\end{enumerate}
\par
\end{proposition}

\begin{proof}
The four main constructions exemplify the four rules.  Cartier realization
uses homogeneous Siu division in the form
\[
                     \chi'=\log|s_F|+\chi_B.
\]
The smooth proof begins with the ideal identity
\[
 J_m\cO_{\mathcal Y}
     =\mathfrak a_m\cO_{\mathcal Y}(mD).
\]
Finite descent uses coefficients of a universal determinant norm, whose
zero polynomial records a common zero in a fibre.  Characteristic-zero
normalization uses a two-sided inequality for nonnegative weighted ideal
values.  None of these steps subtracts two infinite values.
\end{proof}

\begin{remark}[Why divisorial computation is still enough]
\label{rem:ch37-divisorial-computation}
The construction of a candidate is usually checked first on
$X^{\divis}$, where divisor values are finite.  The candidate is then
known independently to be psh, and comparison determines its values at
proper-support points.  Thus divisorial density is not being used to impose
an arbitrary extension by hand.
\end{remark}

\section{Unibranchness and the node--cusp test}
\label{sec:ch37-node-cusp}

Let
\[
                         \nu:X^\nu\longrightarrow X
\tag{37.5}\label{eq:bj-normalization-preview}
\]
be the normalization.  Since $X$ is of finite type over a field, $\nu$
is finite.  If $X$ is unibranch, then there is a unique point of
$X^\nu$ over each point of $X$.  Over algebraically closed $k$, the
finite birational map is radicial, hence a universal homeomorphism.  It
therefore induces a homeomorphism
\[
             \nu^{\An}:(X^\nu)^{\An}\xrightarrow{\ \sim\ }X^{\An}.
\tag{37.6}\label{eq:bj-normalization-analytic-homeomorphism}
\]

The hypothesis is sharp for this reflection mechanism.  If a normalization
has two points over one singular point, then a branch-asymmetric function
upstairs cannot be a pullback from downstairs.  We do not infer from this
obstruction alone a converse to \Cref{thm:bj-final-envelope}.

\begin{example}[Node versus cusp]
\label{ex:bj-node-cusp}
Assume, for the displayed plane equations, that
$\operatorname{char}k\ne2,3$.
\begin{enumerate}
\item An integral nodal cubic has two points $p_+$ and $p_-$ in its
normalization above the node $p$.  The two proper-support endpoints
upstairs are identified downstairs.  Every pulled-back function takes the
same value at $p_+$ and $p_-$, while a GF system on the normalization
can distinguish them.  Full normalization reflection therefore fails.
\item The cuspidal cubic $y^2z=x^3$ has one point in its normalization
above the cusp.  Its completed local inclusion is
\[
                         k[[t^2,t^3]]\subset k[[t]].
\]
The curve is nonnormal but unibranch, and its normalization is a finite
birational universal homeomorphism.  It lies within the theorem.
\end{enumerate}
\end{example}

\begin{remark}[Characteristic of the picture]
\label{rem:ch37-node-cusp-characteristic}
In characteristics $2$ and $3$, use standard local models of an ordinary
node and a unibranch cusp.  The proof uses the number of normalization
branches and the universal-homeomorphism criterion, not the displayed
complex-analytic equations.
\end{remark}


\section{Testing the hypotheses}
\label{sec:ch37-testing}

The theorem combines a geometric condition, two admissibility conditions on
the family, and the ambient scope of Part~IV.  Unibranch sharpness in an ample
class was settled in \Cref{thm:ep-implies-unibranch}, while
\Cref{sec:ch37-node-cusp} isolated the associated normalization obstruction.
The two examples below test the family-level conditions.  No sharpness claim
is made here for the ambient assumptions on the field or the variety.  The
examples are small and are meant to be checked by hand.

\begin{example}[Why a common upper bound is essential]
\label{ex:ch37-pointwise-not-uniform}
Let $X=\mathbf P^1$ and $\theta=c_1(\cO(1))$, so $\deg\theta=1$.  For
$j\in\mathbf Z_{>0}$, take the Fubini--Study function
$\rho_j:=\phi_c$ of \eqref{eq:p1-fs} with $c=j$.  Thus
$\rho_j=0$ off the ray through $\ord_0$ and
$\rho_j(t\ord_0)=\max\{-t,-j\}$; its psh property and all these values
were checked in \Cref{ex:p1-fs-function}.  Put $\psi_j:=\rho_j+j$.  Then
\[
 \psi_j(t\ord_0)=\max\{j-t,\,0\},
 \qquad
 \psi_j\equiv j\ \text{elsewhere}.
\]
Each $\psi_j$ is $\theta$-psh, since adding a constant preserves the
class by \Cref{thm:bj-basic-closure}\,(1).  Each individual $\psi_j$ is
finite at each point, and the family is unbounded above at every point of
$X^{\An}$ but one: $\psi_j\equiv j$ off the ray, and
$\psi_j(t\ord_0)=\max\{j-t,0\}\to+\infty$ for each fixed $t<\infty$,
whereas at the endpoint $v_{0,\triv}=\lim_{t\to\infty}t\ord_0$ one has
$\psi_j(v_{0,\triv})=0$ for \emph{every} $j$.  So $\sup_j\psi_j$ is
$+\infty$ on $X^{\An}\setminus\{v_{0,\triv}\}$ and $0$ at
$v_{0,\triv}$; it is not a psh function, and its usc regularization is
identically $+\infty$, because every neighbourhood of $v_{0,\triv}$
meets the ray at finite parameter.  The
example therefore shows that the boundedness hypothesis cannot simply be
dropped; it does \emph{not} exhibit a pointwise-bounded family that fails to
be uniformly bounded, and by
\Cref{rem:ch37-pointwise-equals-uniform} no such family exists.

The example is deliberately crude: finiteness of every \emph{individual}
member at a given point says nothing about the supremum of the family there.
By \Cref{rem:ch37-pointwise-equals-uniform}, an actual pointwise upper bound
for the whole family would already give the common constant required by
\Cref{def:bj-bounded-family}.  Normalization and boundedness should also not
be conflated.  Compactness arguments remove the additive freedom by passing
to $\PSH_X^{\sup}(\theta)$ in \eqref{eq:normalized-psh-space}; for the
envelope of a prescribed family one imposes a common upper bound, since
normalizing its members separately would change their supremum.
\end{example}

\begin{example}[Nonemptiness is not automatic]
\label{ex:ch37-empty-class}
If $\theta\notin\operatorname{Psef}(X)$, then
$\PSH_X(\theta)=\varnothing$ by
\Cref{lem:ch34-psh-nonempty}, and there is no family to which the theorem
could apply.  The hypothesis that the family be nonempty is thus doing two
things at once: it excludes the empty family, whose supremum is the constant
$-\infty$, and, implicitly, it forces $\theta$ to be pseudoeffective.
It does not force nefness; that stronger conclusion follows only from the
existence of a psh member bounded on both sides.  This is why
\Cref{thm:bj-final-envelope} does not need a positivity hypothesis on
$\theta$ as a separate item: every nonvacuous instance is automatically
pseudoeffective, while nefness is neither assumed nor implied by
nonemptiness alone.
\end{example}

\begin{remark}[Real classes are not a technicality]
\label{rem:ch37-real-classes}
\Cref{thm:bj-final-envelope} is stated for
$\theta\in N^1(X)_{\mathbf R}$, and the real case is not formally reducible
to the rational one by an unoriented limiting argument.  The reduction in
\Cref{sec:ch40-real-to-rational} approximates $\theta$ by rational big
classes $\theta_j$ satisfying the decisive order condition
$\theta_j-\theta\in\Amp(X)$, and then uses closedness of the psh classes
to return to $\theta$.  Inside the ideal-theoretic construction, the
finite stages carry the exact classes $c_1(L)+m^{-1}c_1(A)$, not merely
unspecified $O(m^{-1})$ perturbations.  Thus the real-class theorem rests
on ordered class drift and decreasing-limit closure, not on an assumed
continuity of the envelope operator.
\end{remark}

\begin{remark}[What the theorem does not say]
\label{rem:ch37-not-said}
Three misreadings are worth forestalling.

First, the theorem does not assert that $\sup_i\varphi_i$ is itself psh; it
asserts this of the regularization, and
\Cref{ex:ch36-sup-not-usc} shows the two differ in general.

Second, it does not assert that the supremum is attained, nor that the
regularization agrees with the supremum off a small set in any quantitative
sense.  It is a closure statement, not a regularity statement.

Third, it does not assert that non-unibranch varieties fail every envelope
statement.  What \Cref{thm:ep-implies-unibranch} gives is a converse in the
\emph{ample} case; for a general class on a non-unibranch variety the
question is not addressed here, and the descent used in the proof of
\Cref{thm:bj-final-envelope} simply becomes unavailable.
\end{remark}

\section{Hypothesis ledger}
\label{sec:ch37-hypothesis-ledger}

\begin{proposition}\label{prop:bj-final-hypothesis-ledger}
In \Cref{thm:bj-final-envelope}, the hypotheses have the following exact
scope.
\begin{enumerate}
\item The ground field is algebraically closed and carries the trivial
      absolute value.
\item The variety is integral, projective, and unibranch.  It need not be
      normal or smooth.
\item The numerical class is an arbitrary element of
      $N^1(X)_{\mathbf R}$.  No rationality, bigness, or nefness is
      assumed.  The existence of the displayed nonempty family forces
      $\theta$ to be pseudoeffective by
      \Cref{lem:ch34-psh-nonempty}; for a non-pseudoeffective class the
      envelope-property statement is vacuous.
\item The family may be uncountable.  It is required only to be nonempty
      and bounded above; no monotonicity or common finite locus is assumed.
\item The conclusion concerns the usc regularization of the pointwise
      supremum on the whole analytic space, including proper-support
      points where the value may be $-\infty$.
\end{enumerate}
\par
\end{proposition}

\begin{proof}
These are precisely the assumptions and conclusion in
\Cref{thm:bj-final-envelope}.  If the psh space is empty, there is no
nonempty family and the assertion is vacuous.  No additional case is
created by a non-pseudoeffective class.
\end{proof}

\section{Proof map}
\label{sec:ch37-proof-map}

The order below prevents circularity.

\begin{enumerate}
\item \emph{Real Cartier realization.}
\Cref{chap:bj-cartier-realization} realizes every anti-effective real Cartier
b-divisor $B$ with
nef total class $\pi^*\theta+B_Y$ by a homogeneous psh function.  This
uses no envelope property.
\item \emph{Homogeneous reflection.}
\Cref{chap:bj-reflection} applies the realization theorem to a decreasing homogeneous GF
regularization on a birational model.  An exact drift-cancellation identity
keeps all approximants in one fixed downstairs class.
\item \emph{Carrier reflection.}
The homogeneous Legendre slices of an ordinary psh function are reflected
one by one.  A homogeneous carrier controls the omitted tail; a summable
lower-grid construction preserves decreasing order.
\item \emph{Smooth arbitrary-characteristic theorem.}
\Cref{chap:bj-smooth-characteristic} recalls the multiplier-ideal proof in characteristic zero and
uses asymptotic test ideals on $X\times\mathbf A^1$ in characteristic
$p>0$.  The stages have exact classes $L+A/m$.
\item \emph{Finite normal descent.}
\Cref{chap:bj-finite-descent} proves analytic openness, flattens a finite algebra over a
birational model, and applies determinant norms of universal weighted
section systems.
\item \emph{Normal targets.}
\Cref{chap:bj-alterations} factors a smooth alteration into a projective birational map
followed by a finite normal map.  Birational and finite descent then prove
the envelope property for normal projective varieties.
\item \emph{Universal-homeomorphic normalization.}
\Cref{chap:bj-universal-homeomorphisms,chap:bj-subintegral-normalization}
reflect psh functions through
\eqref{eq:bj-normalization-preview}: Frobenius powers treat positive
characteristic, while the Roberts--Singh additive-power sandwich treats
characteristic zero.
\item \emph{Assembly.}
\Cref{chap:bj-envelope-consequences} pulls the envelope to $X^\nu$, applies normal envelope property,
and reflects the result through
\eqref{eq:bj-normalization-analytic-homeomorphism}.
\end{enumerate}

\paragraph{Proof checkpoint.}
The proof uses three distinct topological facts.  Compactness bounds each
individual psh function above.  Openness of a finite analytic map commutes
pullback with upper semicontinuous regularization.  Homeomorphy of the
unibranch normalization identifies the final envelope pointwise.

\paragraph{Proof checkpoint.}
Arbitrary characteristic enters twice.  Test ideals replace multiplier
ideals on a smooth positive-characteristic source.  The normalization step
then splits into Frobenius and characteristic-zero subintegral arguments.
The determinant norm itself is characteristic-free.

\section{Source note}

The sufficiency theorem and its dependency order are those of the
unpublished manuscript supplied for this project, \emph{The envelope
property for unibranch projective varieties}; the provenance policy is
recorded in \Cref{app:provenance}.  No published-source status is claimed
for that proof.  The definition of the envelope property, the
comparison interface, the increasing-net equivalence, the ample-class
necessity of unibranchness, and the historical conjecture are respectively
\cite[Definition~5.8, Theorem~4.22, Theorem~5.11, Theorem~5.12, and
Conjecture~5.14]{BJ22}.  The published smooth partial result later extended
by the supplied theorem is \cite[Theorem~A]{BJ22b}.  Thus the two directions have
different provenance: the supplied manuscript proves unibranch
sufficiency, while the converse used here only in an ample class is the
published theorem \cite[Theorem~5.12]{BJ22}.  The node--cusp comparison
only illustrates the normalization mechanism and is not substituted for
that converse.
